\sectioncentered*{Введение}
\addcontentsline{toc}{section}{Введение}

% 1. Актуальность
% Требования:
% - обосновать актуальность темы интерфейса интеллектуальной системы;
% - указать контекст: предметная область, тип задач, связь с предыдущими семестрами.
%
% 2. Цель и задачи курсового проекта
% Требования:
% - сформулировать одну цель курсового проектирования;
% - задать 3–6 задач, соответствующих основным разделам пояснительной записки.
%
% 3. Объект, предмет и границы разработки
% Требования:
% - явно указать объект (интеллектуальная система/подсистема);
% - указать предмет (интерфейс, сценарии взаимодействия, архитектура интерфейса);
% - очертить границы реализуемого функционала и прототипа.
%
% 4. Структура пояснительной записки
% Требования:
% - кратко описать содержимое разделов 1–3 и заключения;
% - не дублировать дословно формулировки последующих разделов.

Современный этап развития информационных технологий характеризуется фундаментальным сдвигом парадигмы человеко-машинного взаимодействия. Традиционные графические пользовательские интерфейсы (GUI), основанные на визуальных метафорах и прямом манипулировании элементами управления, постепенно дополняются, а в ряде случаев и замещаются интерфейсами на естественном языке. Особую актуальность данная тенденция приобретает в сфере интерактивных виртуальных сред и компьютерных игр, где сложность игровых механик и объем поступающей информации часто превышают когнитивные способности пользователя по их быстрой обработке.
Актуальность темы курсового проекта обусловлена необходимостью создания новых подходов к проектированию пользовательских интерфейсов, способных не просто отображать статические данные, но и интерпретировать контекст происходящего, предоставляя пользователю релевантную обратную связь в интуитивно понятной форме. Разработка интеллектуальных ассистентов, интегрированных непосредственно в процесс взаимодействия с программным обеспечением, позволяет существенно снизить порог вхождения для новых пользователей и обогатить опыт опытных операторов системы. В условиях роста популярности больших языковых моделей и технологий синтеза речи открываются возможности для создания адаптивных интерфейсов, имитирующих социальное взаимодействие и обладающих элементами эмоционального интеллекта.
Анализ достижений в предметной области показывает, что существующие решения в сфере голосовых помощников (Siri, Google Assistant, Алиса) преимущественно ориентированы на выполнение атомарных команд и поиск информации по прямому запросу пользователя. Однако они обладают существенным ограничением — отсутствием глубокой интеграции с контекстом текущей деятельности пользователя на экране (так называемая проблема контекстной слепоты). В игровой индустрии интерфейсы традиционно представлены в виде HUD (Heads-Up Display) и текстовых чатов, которые являются пассивными средствами отображения информации. Попытки внедрения NPC (неигровых персонажей) с искусственным интеллектом предпринимаются, однако они часто ограничены заранее прописанными сценариями поведения и отсутствием гибкости в диалоге. Проблема создания мультимодального интерфейса, который объединял бы в себе событийную осведомленность, визуальную персонализацию и голосовое взаимодействие в реальном времени, остается открытой инженерной и дизайнерской задачей.
В качестве предметной области для разработки выбрана среда Minecraft, представляющая собой сложную систему с открытым миром и множеством неструктурированных событий, что делает ее идеальным полигоном для тестирования адаптивных интерфейсов.
Целью курсового проекта является проектирование и разработка интеллектуального мультимодального интерфейса ассистента для виртуальной среды, обеспечивающего контекстное аудиовизуальное сопровождение пользователя на основе анализа событийного потока в реальном времени с использованием больших языковых моделей.
Для достижения поставленной цели в ходе выполнения работы необходимо решить следующие задачи:
\begin{itemize}

\item провести анализ предметной области и существующих аналогов интеллектуальных ассистентов в игровых и прикладных средах, выявить недостатки текущих интерфейсных решений и сформулировать требования к проектируемой системе.
\item разработать архитектуру программно-аппаратного комплекса, включающую модули сбора данных из приложения, серверную часть для агрегации событий и принятия решений, а также клиентские интерфейсы управления и отображения.
\item спроектировать и реализовать серверную логику обработки естественного языка, обеспечивающую фильтрацию событий, генерацию контекстно-зависимых реплик и управление эмоциональной окраской синтезируемой речи.
\item разработать пользовательские интерфейсы системы: веб-панель для настройки параметров личности ассистента и мониторинга его работы, а также внутриигровой оверлей для визуализации аватара и субтитров.
\item провести тестирование разработанной системы, оценить задержки передачи данных и реакции интерфейса, а также проанализировать удобство взаимодействия пользователя с разработанным интеллектуальным ассистентом.

\end{itemize}